\chapter{الطاقة الميكانيكية وانحفاظها}\label{ch:g11:mechanical-energy}

لا محرك للأفعوانية بعد التلّ الأول: فكل حلقة وكل اندفاعة
سرعة بعده تُنفَق من الحساب الذي فُتح في الصعود.
ويحوّل هذا الفصل المحاسبة إلى عددين --- \omterm{def:g11:mechanical-energy:kinetic}{الطاقة الحركية} \omterm{def:g10:energy-conservation:energy}{والطاقة}
الكامنة --- وقاعدة واحدة: ما دام \omterm{def:g10:universal-gravitation:weight}{الوزن} وحده يشتغل، فمجموعهما
لا يتغير. \omterm{def:g10:inertia:inventory}{والاحتكاك} هو الرسم؛ وسنتعلم قراءته من
الدفاتر.

\section{الطاقة الحركية}

\begin{definition}[الطاقة الحركية]\label{def:g11:mechanical-energy:kinetic}
جسم كتلته $m$ (\unit{kg}) يتحرك بسرعة $v$ (\unit{m/s}) يحمل
\emph{الطاقة الحركية}\index{طاقة حركية}
\[
E_k = \tfrac12\, m v^2 .
\]
وتحقّق الوحدات: $\unit{kg.m^2/s^2} = \unit{N.m} = \unit{J}$ --- أي \omterm{def:g11:work-of-force:work}{الجول}،
كما في كل \omterm{def:g10:energy-conservation:energy}{طاقة} (\cref{ch:g10:energy-conservation}).
\end{definition}

\begin{example}[رتب القدر]\label{ex:g11:mechanical-energy:orders}
\begin{itemize}
  \item ماشٍ كتلته \qty{70}{kg} بسرعة \qty{1.4}{m/s}:
        $E_k = \tfrac12 \times 70 \times 1.4^2 \approx \qty{69}{J}$؛
  \item سيارة كتلتها \qty{1300}{kg} بسرعة \qty{130}{km/h} $= \qty{36.1}{m/s}$:
        $E_k \approx \qty{8.5e5}{J}$؛
  \item قطار سريع كتلته $m = \qty{3.85e5}{kg}$، بسرعة \qty{320}{km/h}
        $= \qty{88.9}{m/s}$: $E_k \approx \qty{1.5e9}{J}$؛
  \item رصاصة بندقية كتلتها \qty{8.0}{g} بسرعة \qty{800}{m/s}: $E_k \approx \qty{2.6}{kJ}$.
\end{itemize}
فسبع رتب من القدر تفصل النزهة عن القطار. وينمو $E_k$
مثل \emph{مربع} السرعة: فعند \qty{130}{km/h} تحمل السيارة
$4$ ضعف \omterm{def:g10:energy-conservation:energy}{الطاقة} التي لها عند \qty{65}{km/h} --- أي أربعة أضعاف
ما يجب أن تنزعه مكابحها.
\end{example}

\section{مبرهنة الطاقة الحركية}

\begin{theorem}[مبرهنة الطاقة الحركية]\label{thm:g11:mechanical-energy:ketheorem}
\index{مبرهنة الطاقة الحركية}
في جسم كتلته $m$ يتحرك على مستقيم من $A$ إلى $B$،
يساوي تغيّر \omterm{def:g11:mechanical-energy:kinetic}{الطاقة الحركية} الشغلَ الكلي
(\cref{ch:g11:work-of-force}) المتلقّى:
\[
\Delta E_k = E_k(B) - E_k(A) = \sum W_{AB}(\vect F),
\]
والمجموع على كل القوى المؤثرة في الجسم. وتمتدّ المبرهنة إلى
الحركة على أي \omterm{def:g10:relative-motion:trajectory}{مسار}، وهو ما نقبله هنا أيضًا.
\end{theorem}
\admitted

\begin{example}[السقوط الحر، متحقَّقًا منه صورةً صورة]\label{ex:g11:mechanical-energy:freefall}
تُسقَط كرة كتلتها \qty{0.200}{kg} وتُصوَّر؛ وتعطي الصور:
\begin{center}\small
\begin{tabular}{lcccc}
$t$ (\unit{s}) & 0.10 & 0.20 & 0.30 & 0.40\\
الارتفاع المقطوع $h$ (\unit{m}) & 0.049 & 0.196 & 0.441 & 0.785\\
السرعة $v$ (\unit{m/s}) & 0.98 & 1.96 & 2.94 & 3.92\\
$E_k = \tfrac12 mv^2$ (\unit{J}) & 0.096 & 0.384 & 0.864 & 1.54\\
$W = mgh$ (\unit{J}) & 0.096 & 0.385 & 0.865 & 1.54
\end{tabular}
\end{center}
وانطلاقًا من $E_k = 0$، يوافق $E_k$ في كل صورة شغلَ $mgh$
\omterm{def:g10:universal-gravitation:weight}{الوزن}: أي $\Delta E_k = W(\vect P)$، أي\ $v^2 = 2gh$ --- وهي المبرهنة،
متحقَّقًا منها في حدود دقة القراءات.
\end{example}

\begin{remark}[أين يقيم البرهان]\label{rem:g11:mechanical-energy:proof}
تنتج الصيغة العامة --- الصحيحة على أي \omterm{def:g10:relative-motion:trajectory}{مسار}، لا على مستقيم
وحده --- عن قانون نيوتن الثاني وقليل من التحليل، وكلاهما
لاحق في هذا الكتاب (\cref{ch:g12:newtons-laws,%
ch:g12:work-and-energy})؛ وإلى ذلك الحين يكون تحقّق السقوط الحر
ضماننا، ونقبل صيغة \omterm{def:g10:relative-motion:trajectory}{المسار} المنحني ونستعملها بحرية.
\end{remark}

\section{الطاقة الكامنة الثقالية}

\begin{definition}[الطاقة الكامنة الثقالية]\label{def:g11:mechanical-energy:potential}
جسم كتلته $m$ على الارتفاع $z$، مقيسًا إلى أعلى من
\emph{مستوى مرجعي}\index{مستوى مرجعي} مختار حيث $z = 0$، يخزّن
\emph{الطاقة الكامنة الثقالية}\index{طاقة كامنة!ثقالية}
\[
E_p = m g z, \qquad g = \qty{9.81}{N/kg}.
\]
\end{definition}

\begin{remark}[الفروق وحدها تهمّ]\label{rem:g11:mechanical-energy:reference}
الانتقال من $A$ إلى $B$ يغيّر $E_p$ بمقدار $mg(z_B - z_A) =
-W_{AB}(\vect P)$: أي عكس شغل \omterm{def:g10:universal-gravitation:weight}{الوزن}
(\cref{ch:g11:work-of-force}). \omterm{def:g11:mechanical-energy:potential}{والمستوى المرجعي} اختيار حرّ ---
الأرضية، أو مستوى البحر، أو سطح طاولة --- وتغييره يزيح $E_p$ في كل موضع بثابت
يسقط من كل فرق. فضعه حيث تكون الأعداد
أبسط، ولا تحرّكه في منتصف المسألة.
\end{remark}

\begin{example}[كتاب على رفّ]\label{ex:g11:mechanical-energy:shelf}
كتاب كتلته \qty{1.2}{kg} يقع على ارتفاع \qty{1.8}{m} فوق الأرضية، وهي نفسها
على ارتفاع \qty{9.0}{m} فوق الشارع. فبالمرجع عند الأرضية:
$E_p = 1.2 \times 9.81 \times 1.8 \approx \qty{21}{J}$؛ وبالمرجع في الشارع:
$E_p \approx \qty{127}{J}$. والسقوط إلى الأرضية يحرّر \qty{21}{J} في
الحسابين معًا.
\end{example}

\section{الطاقة الميكانيكية وانحفاظها}

\begin{definition}[الطاقة الميكانيكية]\label{def:g11:mechanical-energy:em}
\emph{الطاقة الميكانيكية}\index{طاقة ميكانيكية} لجسم هي
\[
E_m = E_k + E_p = \tfrac12 m v^2 + m g z .
\]
\end{definition}

\begin{theorem}[انحفاظ الطاقة الميكانيكية]\label{thm:g11:mechanical-energy:conservation}
\index{انحفاظ الطاقة الميكانيكية}
إذا كانت \omterm{def:g10:inertia:force}{القوة} الوحيدة التي تشتغل على جسم هي وزنه (بلا \omterm{def:g10:inertia:inventory}{احتكاك}؛ والقوى
الأخرى، كرد فعل \omterm{def:g10:relative-motion:trajectory}{مسار} عديم \omterm{def:g10:inertia:inventory}{الاحتكاك}، عمودية على
الحركة)، انحفظت $E_m$: فمن أجل أي نقطتين $A$ و$B$ من الحركة،
\[
\tfrac12 m v_A^2 + m g z_A = \tfrac12 m v_B^2 + m g z_B .
\]
\end{theorem}

\begin{proof}
بمبرهنة \omterm{def:g11:mechanical-energy:kinetic}{الطاقة الحركية}، $E_k(B) - E_k(A) = W_{AB}(\vect P) =
mg(z_A - z_B) = E_p(A) - E_p(B)$: فما تكسبه $E_k$ تفقده $E_p$.
\end{proof}

\begin{example}[زحلوقة مائية]\label{ex:g11:mechanical-energy:slide}
يبدأ طفل ساكنًا في أعلى زحلوقة عديمة \omterm{def:g10:inertia:inventory}{الاحتكاك} ارتفاعها
$h = \qty{3.2}{m}$. وبالمرجع عند الأسفل: $mgh = \tfrac12 mv^2$، ومنه
$v = \sqrt{2gh} \approx \qty{7.9}{m/s}$ (\qty{29}{km/h}) --- مهما تكن
الكتلة \emph{ومهما يكن شكل الزحلوقة}: مستقيمة أو منحنية
أو حلزونية، فلا يدخل إلا هبوط الارتفاع.
\end{example}

\begin{example}[النوّاس]\label{ex:g11:mechanical-energy:pendulum}
تُسحب كرة نوّاس جانبًا حتى ترتفع $h = \qty{12}{cm}$ فوق
أخفض نقطة لها، ثم تُطلق. وعند الطرفين $v = 0$؛ وعند الأسفل
تكون $mgh$ كلها حركية، $v = \sqrt{2gh} \approx \qty{1.5}{m/s}$؛
وعلى الجانب الآخر تصعد الكرة عائدةً إلى \qty{12}{cm} بالضبط.
\end{example}

\begin{omfigure}
\begin{tikzpicture}[scale=0.95, font=\small]
  \coordinate (O) at (0,0); \draw (-0.5,0) -- (0.5,0);
  \fill (O) circle (1.5pt);
  \draw[dashed, omExo] (230:2.2) arc (230:310:2.2);
  \draw (O) -- (230:2.2); \draw (O) -- (270:2.2); \draw (O) -- (310:2.2);
  \fill[omDef] (230:2.2) circle (2.5pt); \fill[omDef] (270:2.2)
    circle (2.5pt); \fill[omDef] (310:2.2) circle (2.5pt);
  \draw[-Stealth, omThm, thick] (270:2.2) -- ++(1.15,0)
    node[below] {$v_{\max}$};
  \node[above left] at (230:2.2) {$v = 0$};
  \node[above right] at (310:2.2) {$v = 0$};
  \draw[dashed, omExo] (230:2.2) -- ++(-0.9,0)
    (270:2.2) -- ++(-2.4,0);
  \draw[Stealth-Stealth] (-2.2,-2.2) -- (-2.2,-1.685)
    node[midway, left] {$h$};
\end{tikzpicture}

{\small يقايض النوّاس $E_p$ مع $E_k$ ثم بالعكس: كلها كامنة
عند الطرفين، وكلها حركية عند الأسفل، والصعود $h$ نفسه مرتين.}
\end{omfigure}

\begin{omfigure}
\begin{tikzpicture}[scale=0.75, font=\small]
  \draw[thick] plot [smooth, tension=0.7] coordinates
    {(0,3.2) (1.2,1.6) (2.4,0.5) (3.6,1.0) (4.6,2.0) (5.6,0.9)};
  \fill (0,3.2) circle (2pt) node[above right] {$A$};
  \fill (2.4,0.5) circle (2pt) node[below] {$B$}
    (4.6,2.0) circle (2pt) node[above] {$C$};
  \draw[omExo] (-0.2,0) -- (6.0,0);
  \fill[omProp!75] (7.0,0) rectangle (7.45,1.6);
  \fill[omProp!75] (7.9,0) rectangle (8.35,0.25)
    (8.8,0) rectangle (9.25,1.0);
  \fill[omDef!75] (7.9,0.25) rectangle (8.35,1.6)
    (8.8,1.0) rectangle (9.25,1.6);
  \draw[omExo] (6.7,0) -- (9.6,0);
  \draw[dashed, omThm] (6.7,1.6) -- (9.6,1.6) node[above left] {$E_m$};
  \node[below] at (7.225,0) {$A$}; \node[below] at (8.125,0) {$B$};
  \node[below] at (9.025,0) {$C$};
  \fill[omProp!75] (6.9,3.0) rectangle +(0.24,0.24);
  \fill[omDef!75] (6.9,2.5) rectangle +(0.24,0.24);
  \node[right] at (7.2,3.12) {$E_p$}; \node[right] at (7.2,2.62) {$E_k$};
\end{tikzpicture}

{\small مخطط أعمدة \omterm{def:g10:energy-conservation:energy}{الطاقة} لأفعوانية عديمة \omterm{def:g10:inertia:inventory}{الاحتكاك}: عند كل موضع
يتراكم عمودا $E_p$ و$E_k$ ليعطيا المجموع $E_m$ نفسه (الخط المتقطع).}
\end{omfigure}

\begin{example}[سرعة الهبوط في قفزة تزلج]\label{ex:g11:mechanical-energy:skijump}
يغادر قافز تزلج منصة الانطلاق بسرعة $v_A = \qty{25}{m/s}$ ويهبط
على انخفاض $\Delta z = \qty{40}{m}$. وبإهمال \omterm{def:g11:circuits-and-power:resistance}{مقاومة} الهواء،
$v_B = \sqrt{v_A^2 + 2g\,\Delta z} = \sqrt{625 + 785} \approx
\qty{38}{m/s}$ --- ولا حاجة إلى معرفة \omterm{def:g10:relative-motion:trajectory}{مسار} الطيران.
\end{example}

\begin{omfigure}
\begin{tikzpicture}[scale=0.75, font=\small]
  \draw[thick] (0,3.8) .. controls (1.2,2.7) and (1.9,2.4) .. (3.1,2.3);
  \draw[thick] (3.1,1.95) .. controls (4.2,1.55) and (5.5,1.0) .. (7.2,0);
  \draw[dashed, omThm, thick] (3.1,2.3) .. controls (4.3,2.25) and
    (5.2,1.7) .. (5.9,0.75);
  \fill (3.1,2.3) circle (2pt) node[above] {$A$}
    (5.9,0.75) circle (2pt) node[above right] {$B$};
  \draw[dashed, omExo] (3.1,2.3) -- (8.2,2.3) (5.9,0.75) -- (8.2,0.75);
  \draw[Stealth-Stealth] (8.0,0.75) -- (8.0,2.3)
    node[midway, right] {$\Delta z$};
\end{tikzpicture}

{\small قفزة التزلج: مهما يكن المنحني المطار من $A$ إلى $B$، لا يدخل
في ميزانية \omterm{def:g10:energy-conservation:energy}{الطاقة} إلا الهبوط $\Delta z$.}
\end{omfigure}

\section{التبديد بالاحتكاك}

\begin{definition}[التبديد]\label{def:g11:mechanical-energy:dissipation}
يقال عن \omterm{def:g10:energy-conservation:energy}{الطاقة} التي تغادر الحساب الميكانيكي إنها
\emph{مبدَّدة}\index{تبديد}: فهي تعود إلى الظهور طاقةً حرارية في
السطحين المتدالكين وفي الهواء. \omterm{def:g10:energy-conservation:energy}{والطاقة} الكلية محفوظة
(\cref{ch:g10:energy-conservation})؛ أما \omterm{def:g11:mechanical-energy:em}{الطاقة الميكانيكية} وحدها فلا.
\end{definition}

\begin{proposition}[العجز يقيس الاحتكاك]\label{prop:g11:mechanical-energy:deficit}
إذا اشتغلت \omterm{def:g10:inertia:force}{قوة} \omterm{def:g10:inertia:inventory}{احتكاك} $\vect f$ هي أيضًا على امتداد الحركة، فإن
\[
\Delta E_m = E_m(B) - E_m(A) = W_{AB}(\vect f) < 0:
\]
أي أن فقدان \omterm{def:g11:mechanical-energy:em}{الطاقة الميكانيكية} يساوي شغل \omterm{def:g10:inertia:inventory}{الاحتكاك}.
\end{proposition}

\begin{proof}
تصير مبرهنة \omterm{def:g11:mechanical-energy:kinetic}{الطاقة الحركية} الآن $\Delta E_k = W_{AB}(\vect P) +
W_{AB}(\vect f) = -\Delta E_p + W_{AB}(\vect f)$؛ فانقل $\Delta E_p$
إلى الطرف الآخر. \omterm{def:g10:inertia:inventory}{والاحتكاك} يعاكس الحركة، ومن ثم يكون شغله سالبًا.
\end{proof}

\begin{example}[تدقيق زحّافة]\label{ex:g11:mechanical-energy:toboggan}
يبدأ طفل وزحّافة ($m = \qty{40}{kg}$) ساكنين، فيهبطان
$h = \qty{5.0}{m}$ على \qty{12}{m} من الثلج، ويصلان بسرعة
\qty{7.0}{m/s} بدل $\sqrt{2gh} \approx \qty{9.9}{m/s}$
في حالة انعدام \omterm{def:g10:inertia:inventory}{الاحتكاك}. وبالمرجع عند الأسفل:
$\Delta E_m = \tfrac12 \times 40 \times 7.0^2 - 40 \times 9.81
\times 5.0 = 980 - 1962 = \qty{-982}{J}$. ومنه $W(\vect f) =
\qty{-982}{J}$ على مدى \qty{12}{m}: أي \omterm{def:g10:inertia:force}{قوة} \omterm{def:g10:inertia:inventory}{احتكاك} متوسطة
$f = 982/12 \approx \qty{82}{N}$، و\qty{982}{J} من الدفء في
الزلاجتين وفي الثلج.
\end{example}

\begin{omfigure}
\begin{tikzpicture}
\begin{axis}[omaxis, width=8.6cm, height=4.6cm,
    xlabel={$t$ (\unit{s})}, ylabel={$E_m$ (\unit{J})},
    xmin=0, xmax=10, ymin=0, ymax=1250, legend pos=south west]
  \addplot[omThm, thick, domain=0:10, samples=2] {1000};
  \addplot[omDef, thick, domain=0:10, samples=60] {1000*exp(-0.18*x)};
  \legend{without friction, with friction}
\end{axis}
\end{tikzpicture}

{\small $E_m$ على امتداد حركة: مستوية بلا \omterm{def:g10:inertia:inventory}{احتكاك}، ومتناقصة
معه --- والهبوط في كل لحظة هو الحرارة المنتَجة حتى تلك اللحظة.}
\end{omfigure}

\begin{method}[محاسبة الطاقة]\label{met:g11:mechanical-energy:bookkeeping}
\begin{enumerate}
  \item اختر مستوًى مرجعيًّا واحفظه؛ وسمِّ $A$ (المعطيات) و$B$
        (المطلوب).
  \item اكتب $E_m = \tfrac12 mv^2 + mgz$ عند $A$ وعند $B$.
  \item بلا \omterm{def:g10:inertia:inventory}{احتكاك}: $E_m(A) = E_m(B)$؛ وحُلّ --- وعادةً
        $v_B = \sqrt{v_A^2 + 2g(z_A - z_B)}$.
  \item ومع \omterm{def:g10:inertia:inventory}{احتكاك}: $E_m(B) - E_m(A) = -f\ell$ (حيث $\ell$ = طول \omterm{def:g10:relative-motion:trajectory}{المسار})؛
        وحُلّ لإيجاد المجهول.
  \item وتحقّق: يتراكم عمودا $E_k$ و$E_p$ ليعطيا $E_m$، مستويًا ما لم يكن ثمة \omterm{def:g10:inertia:inventory}{احتكاك}.
\end{enumerate}
\end{method}

\section{تمارين}

\begin{exercise}[$\star$]\label{exo:g11:mechanical-energy:1}
احسب \omterm{def:g11:mechanical-energy:kinetic}{الطاقة الحركية} لكلٍّ من: (أ) دراجة نارية صغيرة وراكبها كتلتهما \qty{90}{kg} بسرعة
\qty{25}{km/h}؛ (ب) سيارة كتلتها \qty{1200}{kg} بسرعة \qty{130}{km/h}. (ج) وبأي
عامل تتغير $E_k$ حين تتضاعف السرعة؟ وحين تُثلَّث؟
\end{exercise}

\begin{exercise}[$\star$]\label{exo:g11:mechanical-energy:2}
متسلق كتلته \qty{75}{kg} يكسب \qty{850}{m} من الارتفاع. احسب \omterm{def:g10:signals-and-waves:oscilloscope}{كسب}
\omterm{def:g10:energy-conservation:energy}{الطاقة} الكامنة. وهل يتوقف على \omterm{def:g11:mechanical-energy:potential}{المستوى المرجعي}؟ وعلى الطريق؟
\end{exercise}

\begin{exercise}[$\star$]\label{exo:g11:mechanical-energy:3}
يسقط أصيص من ارتفاع \qty{20}{m} من شرفة. بإهمال \omterm{def:g11:circuits-and-power:resistance}{مقاومة} الهواء،
جِد سرعته عند الأرض \omterm{def:g10:orders-of-magnitude:unit}{بوحدة} \unit{m/s} وبوحدة \unit{km/h}. وأين ذهبت الكتلة؟
\end{exercise}

\begin{exercise}[$\star$]\label{exo:g11:mechanical-energy:4}
تكبح سيارة كتلتها \qty{1200}{kg} من \qty{25}{m/s} إلى السكون في \qty{50}{m}.
استعمل مبرهنة \omterm{def:g11:mechanical-energy:kinetic}{الطاقة الحركية} لإيجاد شغل المكابح، ثم
\omterm{def:g10:inertia:force}{قوة} الكبح المتوسطة.
\end{exercise}

\begin{exercise}[$\star$]\label{exo:g11:mechanical-energy:5}
تُطلق كرة نوّاس من ارتفاع \qty{8.0}{cm} فوق أخفض نقطة لها. جِد
سرعتها عند الأسفل والارتفاع المبلوغ على الجانب الآخر.
\end{exercise}

\begin{exercise}[$\star\star$]\label{exo:g11:mechanical-energy:6}
باستعمال \cref{ex:g11:mechanical-energy:orders}: بأي عامل تتجاوز \omterm{def:g11:mechanical-energy:kinetic}{الطاقة الحركية}
للقطار طاقةَ الرصاصة؟ وإلى أي ارتفاع كانت $E_k$ الخاصة به
ستقدر على رفعه ($h = v^2/2g$)؟
\end{exercise}

\begin{exercise}[$\star\star$]\label{exo:g11:mechanical-energy:7}
زحلوقتان عديمتا \omterm{def:g10:inertia:inventory}{الاحتكاك} تهبطان من منصة ارتفاعها \qty{3.0}{m} نفسها، إحداهما
مستقيمة حادّة، والأخرى طويلة متعرّجة. قارن سرعتي الوصول
(واحسبهما) وزمني الوصول (دون حساب).
\end{exercise}

\begin{exercise}[$\star\star$]\label{exo:g11:mechanical-energy:8}
تُقذف كرة شاقوليًّا إلى أعلى بسرعة \qty{12}{m/s}. جِد أقصى ارتفاع
تبلغه، ثم سرعتها عند \qty{4.0}{m} --- صعودًا ونزولًا.
\end{exercise}

\begin{exercise}[$\star\star$]\label{exo:g11:mechanical-energy:9}
يغادر قافز تزلج منصة الانطلاق بسرعة \qty{23}{m/s} ويهبط على انخفاض \qty{35}{m}.
تنبّأ بسرعة الهبوط دون \omterm{def:g11:circuits-and-power:resistance}{مقاومة} الهواء؛ فهل القيمة \omterm{def:g11:lenses-and-eye:images}{الحقيقية}
أكبر أم أصغر، ولماذا؟
\end{exercise}

\begin{exercise}[$\star\star$]\label{exo:g11:mechanical-energy:10}
زحّافة كتلتها \qty{45}{kg} تبدأ ساكنة، وتهبط \qty{6.0}{m} على
\qty{15}{m} من المنحدر، فتصل بسرعة \qty{8.0}{m/s}. جِد \omterm{def:g11:mechanical-energy:em}{الطاقة
الميكانيكية} المفقودة، ثم \omterm{def:g10:inertia:force}{قوة} \omterm{def:g10:inertia:inventory}{الاحتكاك} المتوسطة. وأين ذهبت \omterm{def:g10:energy-conservation:energy}{الطاقة}؟
\end{exercise}

\begin{exercise}[$\star\star$]\label{exo:g11:mechanical-energy:11}
وتد غاليليو: نوّاس يُطلق من ارتفاع \qty{15}{cm} فوق أخفض نقطة له
يلاقي وتدًا يحجز النصف الأعلى من الخيط عند الوضع الشاقولي. جِد
السرعة عند الأسفل والصعود وراء الوتد؛ فما الذي يغيّره الوتد،
وما الذي لا يغيّره؟
\end{exercise}

\begin{exercise}[$\star\star\star$]\label{exo:g11:mechanical-energy:12}
ينزل متزلج على لوح في حوض نصف أسطواني من ارتفاع \qty{2.5}{m}؛ وكل عبور
(من جانب إلى جانب) يبدّد $10\%$ من \omterm{def:g11:mechanical-energy:em}{الطاقة الميكانيكية}. جِد
السرعة عند أسفل النزول الأول، ثم --- بحساب
ارتفاعات الصعود المتتالية --- أول عبور ينتهي دون \qty{1.0}{m}.
\end{exercise}

\begin{exercise}[$\star\star\star$]\label{exo:g11:mechanical-energy:13}
من قمة جرف ارتفاعه \qty{45}{m}، يُقذف حجر بسرعة \qty{15}{m/s} ---
إلى أعلى، أو أفقيًّا، أو بأي زاوية. بيّن أن سرعة الهبوط
واحدة في كل الحالات واحسبها. وما الذي \emph{يتوقف} على الزاوية؟
\end{exercise}

\begin{exercise}[$\star\star\star$]\label{exo:g11:mechanical-energy:14}
متزلج منحدرات كتلته \qty{80}{kg} يبدأ ساكنًا وينزل \qty{120}{m}
من الهبوط على \omterm{def:g10:relative-motion:trajectory}{مسار} طوله \qty{800}{m}، في مواجهة \omterm{def:g10:inertia:force}{قوة} \omterm{def:g10:inertia:inventory}{احتكاك} ثابتة
\qty{65}{N}. جِد سرعة الوصول والجزء المبدَّد من $E_m$.
\end{exercise}

\begin{exercise}[$\star\star\star$]\label{exo:g11:mechanical-energy:15}
ترفع محطة ضخّ وتخزين $V = \qty{2.0e6}{m^3}$ من الماء
(\qty{1000}{kg} لكل \unit{m^3}) عبر \qty{300}{m}. احسب
\omterm{def:g10:energy-conservation:energy}{الطاقة} المخزونة \omterm{def:g11:work-of-force:work}{بالجول} وبوحدة \unit{kW.h}؛ وإذا أُعيد تعنيفها \omterm{def:g10:energy-conservation:efficiency}{بمردود}
$85\%$ (\cref{ch:g10:energy-conservation})، فكم بيتًا يسحب
\qty{10}{kW.h} في اليوم تغذّي مدة يوم واحد؟
\end{exercise}

\section{مسألة: تصميم الحلقة}

\begin{problem}[{مسألة نهاية الأسبوع --- تدقيق طاقي لأفعوانية:
ارتفاع الإطلاق الذي تمليه الحلقة، والسرعات التي يَعِد بها، والاحتكاك الذي
يعترف به حسّاسان، والمكابح التي تغلق الحساب ---
والركوب كله معلَّق بعدد واحد}]
\label{pb:g11:mechanical-energy:1}
يُطلق قطار كتلته $m = \qty{2.0e3}{kg}$ ساكنًا من تلّ إطلاق
ارتفاعه $H$ فينساب، بلا محرك، عبر حلقة دائرية عمودية
نصف قطرها $R = \qty{8.0}{m}$ وقمتها على الارتفاع $2R$.
ويوفّر تحليل الحركة الدائرية (السنة القادمة) المعطى الوحيد الذي نستعيره: فعند
قمة الحلقة لا يضغط القطار على \omterm{def:g10:relative-motion:trajectory}{المسار} إلا إذا كان
$v_{\text{القمة}}^2 \geq gR$. ويهمل الجزآن I وII \omterm{def:g10:inertia:inventory}{الاحتكاك}.

\medskip
\noindent\textbf{الجزء I --- التلّ الذي تمليه الحلقة.}
\begin{enumerate}
  \item احسب أدنى سرعة آمنة عند قمة الحلقة.
  \item بالأرض مستوًى مرجعيًّا: احسب \omterm{def:g11:mechanical-energy:em}{الطاقة الميكانيكية} للقطار
        عند قمة الحلقة بتلك السرعة الدنيا.
  \item استنتج $H_{\min}$ وبيّن أن $H_{\min} = 5R/2$ --- فلا كتلة ولا $g$.
  \item يأخذ الفريق $H = \qty{24}{m}$: احسب سرعة قمة الحلقة الفعلية
        والنسبة $v_{\text{القمة}}^2/(gR)$.
  \item وأين ذهبت $m$؟ ولماذا يناسب التصميم القطارات الممتلئة
        والفارغة سواءً؟
\end{enumerate}

\medskip
\noindent\textbf{الجزء II --- بماذا يَعِد الارتفاع} ($H = \qty{24}{m}$).
\begin{enumerate}[resume]
  \item ما السرعة عند أسفل الحلقة ($z = 0$)، \omterm{def:g10:orders-of-magnitude:unit}{بوحدة} \unit{m/s} و
        \omterm{def:g10:orders-of-magnitude:unit}{بوحدة} \unit{km/h}؟
  \item وما السرعة عند جانب الحلقة ($z = R$)؟
  \item احسب $E_p$ و$E_k$ عند الإطلاق وعند الأسفل والجانب والقمة
        في الحلقة؛ وتحقّق من أن كل زوج يجمع إلى $mgH$.
  \item ويلي الحلقةَ نتوء ارتفاعه \qty{12}{m}: فما السرعة عند قمته؟
  \item أعطِ الصيغة العامة للمقدار $v(z)$؛ وأين يكون الركوب
        أسرع، وأين أبطأ؟
\end{enumerate}

\medskip
\noindent\textbf{الجزء III --- تدقيق \omterm{def:g10:inertia:inventory}{الاحتكاك}.}
حسّاسان على الارتفاع نفسه $z = \qty{2.0}{m}$، يفصل بينهما \qty{60}{m} من
\omterm{def:g10:relative-motion:trajectory}{المسار}، يقرآن $v_1 = \qty{20.4}{m/s}$ ثم $v_2 = \qty{19.2}{m/s}$.
\begin{enumerate}[resume]
  \item احسب تغير \omterm{def:g11:mechanical-energy:em}{الطاقة الميكانيكية} بين الحسّاسين.
        وما وجه الذكاء في وضعهما على الارتفاع نفسه؟
  \item استنتج شغل \omterm{def:g10:inertia:inventory}{الاحتكاك} \omterm{def:g10:inertia:force}{وقوة} \omterm{def:g10:inertia:inventory}{الاحتكاك} المتوسطة.
  \item وأين ذهبت \omterm{def:g10:energy-conservation:energy}{الطاقة} المفقودة؟ سمِّ الشكل والأماكن.
  \item وأيّ جزء من $E_m$ الخاصة بالقطار يضيع في كل \qty{60}{m}؟ ولماذا
        يكون $H_{\min}$ العاري غير آمن في العالم الحقيقي المحتكّ؟
  \item يفصل نحو \qty{90}{m} من \omterm{def:g10:relative-motion:trajectory}{المسار} بين الإطلاق وقمة الحلقة: قدّر
        $E_m$ هناك؛ فهل ينقذ هامش السؤال 4 الحلقةَ بعدُ؟
\end{enumerate}

\medskip
\noindent\textbf{الجزء IV --- إغلاق الحساب.}
الحسّاس الأخير، عند مدخل \omterm{def:g10:relative-motion:trajectory}{مسار} الكبح الأفقي المستقيم
($z = 0$)، يقرأ \qty{18.0}{m/s}.
\begin{enumerate}[resume]
  \item احسب \omterm{def:g11:mechanical-energy:kinetic}{الطاقة الحركية} للقطار هناك.
  \item ويجب أن توقفه المكابح في $d = \qty{45}{m}$: احسب
        \omterm{def:g10:inertia:force}{قوة} الكبح (الثابتة) المطلوبة.
  \item تحقّق من الراحة: احسب التباطؤ $a = F/m$؛ وقارنه بالمقدار $g$.
  \item وكم \omterm{def:g10:energy-conservation:energy}{طاقة} بدّدها \omterm{def:g10:inertia:inventory}{احتكاك} \omterm{def:g10:relative-motion:trajectory}{المسار} على مدى الركوب
        كله؟ وتحقّق من أن \omterm{def:g10:inertia:inventory}{الاحتكاك} $+$ المكابح $= mgH$، حتى \omterm{def:g11:work-of-force:work}{الجول}.
  \item بيت القصيد، في جملة واحدة: أيّ عدد واحد ثبّت أمان
        الحلقة وكل السرعات وفاتورة \omterm{def:g10:inertia:inventory}{الاحتكاك} \omterm{def:g10:inertia:force}{وقوة} الكبح ---
        وما ثمن إهمال \omterm{def:g10:inertia:inventory}{الاحتكاك}؟
\end{enumerate}
\end{problem}
